\subsection{Representations in a separated quasi-complete space}

PROPOSITION 7. --- Let \(\varrho\) be a continuous representation of G in a separated, locally convex, and quasi-complete space E. The sum of the finite-dimensional subrepresentations of E is dense in E.

Let \(x \in \mathrm{E}\) and let U be an open neighborhood of \(x\) . The set of measures \(\nu \in \mathcal{M}(\mathrm{G})\) such that \(\varrho(\nu)x \in \mathrm{U}\) is open in \(\mathcal{M}(\mathrm{G})\) for the topology of compact convergence (cf. no. 2 of V, p. 400). It contains \(\varepsilon_e\) ; hence it contains a measure of the form \(\nu = f_1 \cdot \mu\) where \(f_1 \in \mathcal{C}(\mathrm{G})\) (INT, VIII, p. 171, § 4, n° 7, prop. 19) and, consequently, it contains a measure of the form \(f_2 \cdot \mu\) where \(f_2\) is a G-finite function (prop. 6 of V, p. 464 and prop. 4 of V, p. 462).

The subrepresentation F of \(\gamma_{G}\) generated by \(f_{2}\) is finite-dimensional. Let \(\widetilde{F}\) be the image of the linear map \(f \mapsto \varrho(f)x\) from F into E. The space \(\widetilde{F}\) is finite-dimensional, and it contains the element \(\varrho(f_{2})x\) belonging to U. Since \(\varrho(g)\varrho(f) = \varrho(\gamma_{\mathrm{G}}(g)f)\) for all \(g \in G\) and every \(f \in F\) (formula (1) of V, p. 406), the space \(\widetilde{F}\) is a subrepresentation of \(\varrho\) which intersects U. The proposition is proved.

COROLLARY 1. --- Let π be a continuous irreducible representation of G in a locally convex separated quasi-complete space E. The space E is finite-dimensional.

COROLLARY 2. --- Let \(\varrho\) be a continuous representation of G in a locally convex separated quasi-complete space E, and let \(\pi\) be a continuous irreducible representation of G.

\begin{enumerate}
\def\labelenumi{\alph{enumi})}
\item
  The G-morphism \(p_{\pi} = \dim(\pi)\varrho(\overline{\chi}_{\pi})\) from E into E is a continuous projector of E whose image is the \(\pi\)-isotypic component of \(\varrho\) ;
\item
  If \(\varrho\) is a unitary representation, then the projector \(p_{\pi}\) is the orthogonal projector of E with image \(\mathrm{M}_{\pi}(\varrho)\) .
\end{enumerate}

Let us prove \(a\) ). The linear map \(p_{\pi} = \dim(\pi)\varrho(\overline{\chi}_{\pi})\) is well defined, since E is quasi-complete and G is compact (V, p. 401). It is an element of \(\mathrm{Hom}_{\mathrm{G}}(\varrho,\varrho)\) since the character of \(\pi\) is a continuous central function on G.

Let \(\varpi\) be a finite-dimensional subrepresentation of E and F its space. The map \(p_{\pi}\) induces, by passing to subspaces, the endomorphism \(\dim(\pi)\varpi(\overline{\chi}_{\pi})\) of F. This endomorphism is therefore zero if \(\varpi\) is not isomorphic to \(\pi\) and is the identity map of F otherwise (indeed, this is the case if \(\varpi\) is irreducible by cor. 5 of V, p. 460, and the general case follows from this since \(\varpi\) is semisimple by prop. 1 of V, p. 457).

Finally, since the sum of the finite-dimensional subrepresentations of E is dense in E (prop. 7) and since \(p_{\pi}\) is continuous, we conclude that \(p_{\pi}\) is a projector whose image is the \(\pi\)-isotypic component of \(\varrho\) . If E is a Hilbert space, the projector \(p_{\pi}\) is Hermitian by Cor. 1 of V, p. 458, hence it is an orthogonal projector (Lemma 3, (ii) of I, p. 133).

COROLLARY 3. --- Let \(\varrho\) be a continuous representation of G in a locally convex separated quasi-complete space E. The endomorphism

\[
x \mapsto \int_ {\mathrm{G}} \varrho (g) x d \mu (g)
\]

of E is a projector of E whose image is the space \(E^{G}\) of invariant vectors in E. In particular, if E is finite-dimensional, then

\[
\mathrm{dim} (\mathrm{E} ^ {\mathrm{G}}) = \int_ {\mathrm{G}} \chi_ {\varrho} (g) d \mu (g).
\]

The first assertion is the special case of the preceding corollary when \(\pi\) is the trivial representation of G, of dimension 1. The second follows from it since the dimension of \(E^{G}\) is the trace of the orthogonal projector onto \(E^{G}\) , that is,

\[
\mathrm{dim} (\mathrm{E} ^ {\mathrm{G}}) = \mathrm{Tr} \Bigl (\int_ {\mathrm{G}} \varrho (g) d \mu (g) \Bigr) = \int_ {\mathrm{G}} \chi_ {\varrho} (g) d \mu (g).
\]

In particular, when E is finite-dimensional, there exists a vector \(x \neq 0\) in E such that \(\varrho(g)x = x\) for all \(g \in G\) if and only if

\[
\int_ {\mathrm{G}} \chi_ {\varrho} (g) d \mu (g) \neq 0.
\]

COROLLARY 4. --- Let \(\varrho\) be a unitary representation of G in a Hilbert space E. The space E is the Hilbert sum of the \(\pi\)-isotypic components \(\mathrm{M}_{\pi}(\varrho)\) for \(\pi \in \widehat{\mathrm{G}}\) . In particular, the representation \(\varrho\) is semisimple.

The isotypic components of \(\varrho\) corresponding to non-isomorphic irreducible representations of G are orthogonal (prop. 8 of V, p. 394). Since every finite-dimensional unitary representation of G is semisimple (prop. 1 of V, p. 457), the sum of the isotypic components \(M_{\pi}(\varrho)\) for \(\pi \in \widehat{G}\) is dense in E (prop. 7). The corollary follows from this.
% semantic-fragments: legacy-input-seam
\relax{}
